\documentclass[12pt]{article}
\usepackage[utf8]{inputenc}
\usepackage[T1]{fontenc}
\usepackage{hyperref}
\title{Kauwera Language Resources for K-12 Curriculum Development}
\author{Compiled by Claude Sonnet 5.5 (claude-sonnet-5-5) for Andr\'as Kornai}
\date{}
\begin{document}
\maketitle

\section{Introduction}
Kauwera (ISO 639-3: xau; Glottocode: kauw1242) is a Papuan language of Indonesia in the Greater Kwerba family. Wikipedia cites about 400 speakers (1987). Joshua Project lists a Kauwerawec people group of about 700. Joshua Project also describes the language as unwritten. Alternate names include Kabera, Kaowerawedj, Kauwerawec and Kawera. Searches found very little material, and nothing aimed at schools.

\section{Primary Educational Resources}
No primers, textbooks or school materials in Kauwera turned up in the searches. The language is not known to be taught in schools.

\section{Linguistic and Reference Materials}
\begin{itemize}
\item Overview of the language: \url{https://en.wikipedia.org/wiki/Kauwera_language}
\item Glottolog record (classification, status): \url{https://glottolog.org/resource/languoid/id/kauw1242}
\item Kwerba languages (family context): \url{https://en.wikipedia.org/wiki/Kwerba_languages}
\item Kwerbic languages (the broader Greater Kwerba family): \url{https://en.wikipedia.org/wiki/Kwerbic_languages}
\item New Guinea World, Kwerba (University of Zurich): \url{https://newguineaworld.linguistik.uzh.ch/families/northwest-new-guinea/foja-range/west-foja-range/apauwar-river/kwerba}
\item ASJP word list entry (Kwerba Kauwerawet I): \url{https://asjp.clld.org/languages/KWERBA_KAUWERAWET_I}
\end{itemize}
Wikipedia and Glottolog mention early sources (Le Roux 1926; van Eechoud 1962). No URLs for these were seen.

\section{Religious and Community Resources}
Joshua Project lists the people group and notes that Bible translation is needed: \url{https://joshuaproject.net/languages/xau}. No scripture or community organisation was found.

\section{Digital and Audio Resources}
Wikitongues has a page for the code: \url{https://wikitongues.org/languages/xau/}. I did not confirm whether it holds any recordings. No other audio, video or corpus was found.

\section{Curriculum Development Context}
There is no known written standard, orthography or school programme. Any curriculum would have to begin with documentation, orthography design and community consultation, using the Indonesian national language as the bridge language.
\end{document}